\documentclass[10pt,a4paper]{amsart}
\usepackage[margin=19mm]{geometry}
\usepackage{amsmath,amssymb,mathtools}
\usepackage[scr=rsfs,cal=euler]{mathalfa}
\usepackage{xfrac}
\usepackage[unicode,hidelinks,bookmarks=false]{hyperref}
\newcommand{\ie}{i.e.\ }
\newcommand{\cf}{cf.\ }
\newcommand{\resp}{respectively\ }
\newcommand{\Subsection}{\subsection}
\providecommand{\nobreakdash}{\nobreak\hbox{-}}
\setlength{\emergencystretch}{2em}
\multlinegap0pt
\usepackage{bm}
\usepackage{bbm}
\usepackage{dsfont}
\usepackage{xspace}
\usepackage{cancel}
\usepackage{mathtools}
\usepackage{amsthm}
\makeatletter
\@ifundefined{theorem}{\newtheorem{theorem}{Theorem}}{}
\@ifundefined{corollary}{\newtheorem{corollary}[theorem]{Corollary}}{}
\@ifundefined{lemma}{\newtheorem{lemma}[theorem]{Lemma}}{}
\@ifundefined{proposition}{\newtheorem{proposition}[theorem]{Proposition}}{}
\makeatother
\title[Brownian paths as loop-decorated SLEs]{Brownian paths as loop-decorated SLEs}
\author{Nathana\"el Berestycki, Isao Sauzedde}
\date{Source: arXiv:2602.04673v1 (CC BY 4.0)}
\begin{document}
\maketitle

\noindent\emph{Source and licence.} Brownian paths as loop-decorated SLEs. Version arXiv:2602.04673v1 (CC BY 4.0), distributed under the Creative Commons Attribution 4.0 International licence, as recorded in the arXiv licence metadata for this exact version and shown on the abstract page https://arxiv.org/abs/2602.04673v1. Full rights notice: licenses/arxiv-2602.04673-CC-BY-4.0.txt.\par\smallskip

\noindent\textit{Editorial scope.} Counted unit: the Proposition whose source label is \texttt{prop:vol2} in the article (source0.tex lines 2992--3009, source label \texttt{prop:vol2}), delivered together with the article's own complete original proof of it (source0.tex lines 3010--3147, one \texttt{proof} environment of 138 lines). No mathematics is written, repaired or re-derived: statement and proof are the article's own text line for line. Required context retained uncounted: coro:vol (lines 2953--2959). Every definition, display and label of those lines is retained unchanged. Omitted from this package and not redistributed: 1--2952, 2960--2991, 3148--3657. Nothing inside a retained excerpt is omitted or altered: every retained body is delivered exactly as the article's lines are written, so no omission receipt of any kind accompanies this package. Citations: the retained excerpts carry no citation command at all, so no bibliography entry of the article is transcribed and no reference is redistributed. Reference adaptation: none was needed and none was made; every reference inside the delivered unit resolves inside it, so no unresolved reference or citation is produced. Printed slips kept literal: no printed word, symbol, hypothesis or number of the counted statement or of its proof was altered. Layout and class: the article's own class is \texttt{article} with packages including inputenc, xcolor, mathtools, mathrsfs, bbm, geometry, amsmath, amsfonts, amssymb, graphics, amsthm, comment, hyperref, caption; the wrapper is the lane's pinned \texttt{amsart} editorial wrapper at 10pt on A4, which restates the theorem-like environments the retained text needs and adds the article's own macro definitions that the retained text needs (none), with extra wrapper packages (bm, bbm, dsfont, xspace, cancel, mathtools). The delivered numbering is the unit's own. Assets: no retained span contains a figure environment, a graphics-inclusion command, a PDF-page inclusion, a table or a TikZ picture; no image file is copied into this package, the package declares no asset dependency and no image file is redistributed with it. Licence: version arXiv:2602.04673v1 of this article is distributed under the Creative Commons Attribution 4.0 International licence, as recorded in the arXiv licence metadata for this exact version and shown on the abstract page https://arxiv.org/abs/2602.04673v1. A full rights notice is delivered at licenses/arxiv-2602.04673-CC-BY-4.0.txt. Characters: every retained excerpt is pure ASCII, so the delivered engine, XeLaTeX with its default Latin Modern text font, reports no missing character.
\begin{corollary}
    \label{coro:vol}
    There exists $C',C_1<\infty$ and $\nu>0$ such that for all positive integers $n$ and $k\geq C' \log(n)$,
    \[
\mathbb{E}[\#\{ x\in n D: d(x,\gamma^{nD})\leq k  \}] \leq C_1  k^\nu n^{2-\nu}.
    \]
\end{corollary}

\begin{proposition}
    \label{prop:vol2}
    There exists $C<\infty$ and $\nu>0$ such that for all positive integers $n$ and $k$,
    \[
    \mathbb{E}\Big[ \#   \{   \mathsf{L}\in \mathbb{L}^{nD}_{\gamma^{nD}}: t_\mathsf{L}=k \}  \Big]
    \leq C n^{2-\nu}   k^{-2}
    \max(\sqrt{k},   \log(n))^\nu .
    \]
    As a consequence, (recalling that $T_{\leq \delta} (\mathcal{X}^n)$ denotes the total time spent by $\mathcal{X}^n$ on loops of duration less than $\delta n^2$), there exists $C<\infty$ such that for all $n\geq 1$ and all $\delta>0$,
    \begin{equation}
    \label{eq:temp:Ebound}
    \mathbb{E}[T_{\leq \delta}(\mathcal{X}^n)]
    %=n^{-2}
    %\sum_{k=1}^{\lfloor \delta n^2 \rfloor}  k \mathbb{E}[\# \{  \mathsf{L}\in \mathbb{L}^{nD}_{\gamma^{nD}}: t_{\mathsf{L}}=k  \}]
    %
    \leq C ((\log\log n) (\log n)^\nu n^{-\nu}+ \delta^{\frac{\nu}{2}}).
    \end{equation}
\end{proposition}

\begin{proof}
    We assume $k$ even during the proof, for otherwise the left-hand side is equal to $0$.
    Let $\mathsf{L}$ be distributed as a simple random walk loop on $\mathbb{Z}^2$ started from $0$ with $k$ steps under $\mathbb{P}_{k}$, and let
    \[r_\mathsf{L}\coloneqq \max \{ |\mathsf{L}(t)-\mathsf{L}(0)|:   t\in \{0,\dots, k\} \}.\]
    
    
    The following estimate on the tail of $r_\mathsf{L}$ is classical, and we include it here for completeness. 
    
    
\begin{lemma}\label{L:ex_hoeffding} There exists a universal constant $C>0$ such that 
    \begin{equation} 
    \label{eq:temp:hoef}
    \mathbb{P}_k( r_{\mathsf{L}} \geq R \sqrt{k} ) \leq C \exp(-R^2).
    \end{equation}
\end{lemma}

\begin{proof}[Proof of Lemma \ref{L:ex_hoeffding}]
    Let $(X_n, n\ge 0)$ denote simple random walk on $\mathbb{Z}^2$ starting from $X_0 = 0$, and let $M_n = 
    \max \{|X_j|:0
    \le j \le n \}$.  
    By symmetry and reversibility of a random walk bridge of length $k$, it suffices to prove that 
    \begin{equation}\label{eq:AH}
    \mathbb{P}( M_{k/2} \geq R \sqrt{k} | X_k = 0) \leq (C/2) \exp(-R^2).
    \end{equation}
    (Recall that $k$ is assumed even.) 
    Now, let $j$ denote $k/2$ if $k/2$ is even, and $k/2+1$ otherwise. 
    %Let $\tau = \inf\{n \text{even}: |X_
    %n| \ge R \sqrt{k}\}$. 
    Applying the Markov property, and letting $p_n(x,y)$ denote the transition probability for random walk on $\mathbb{Z}^2$, 
    \begin{align*}
    \mathbb{P}( M_{k/2} \geq R \sqrt{k} ; X_k = 0) 
    &\le \mathbb{P}( M_j 
    \geq  R \sqrt{k} ; X_k = 0) \\
    & \leq \mathbb{E}( 1_{
    \{M_{j} \ge R\sqrt{k}\}} p_{k-j}(X_j,0))  .
    \end{align*}
    Now, using Cauchy--Schwarz and the reversibility of simple random walk on $
    \mathbb{Z}^2$ it is well known that for any $x\in \mathbb{Z}^2$ with the same parity as 0 we have $p_{2n}(x,0) \le p_{2n} (0,0)$. This also trivially holds when $x$ does not have the same parity as 0. Hence (since we have chosen $j$ even, hence $k-j$ is also even), we deduce that 
    \begin{align}\label{eq:tailAH}
    \mathbb{P}( M_{k/2} \geq R \sqrt{k} ; X_k = 0)  &\le p_{k-j} (0,0) \mathbb{P}( M_j 
    \geq  R \sqrt{k} )
    \end{align}
    By the maximal Azuma--Hoeffding inequality,
    $$
    \mathbb{P}( M_j 
    \geq  R \sqrt{k} )\le \exp ( - 2\tfrac{R^2k}{j}) \le \exp ( - R^2).
    $$
    Also, by the local central limit theorem, there exists a universal constant $C>0$ such that $p_{k-j}(0,0) \le C p_k(0,0)$. Plugging into \eqref{eq:tailAH} we get
    \[
    \mathbb{P}( M_{k/2} \geq R \sqrt{k} ; X_k = 0)  \le C \exp(-R^2) p_k(0,0).
    \]
    Dividing by $
    P(X_k =0) = p_k(0,0)$ to get the conditional probability, this proves \eqref{eq:AH} (changing the value of $C$ by a factor two) as desired. 
\end{proof}

    
    Now let us return to the proof of Proposition \ref{prop:vol2}.    For any $k\geq 1$, 
    
    \begin{equation}
    \label{eq:temp:mass}
    \mathbb{E}[\# \{ \mathsf{L}\in \mathbb{L}^{\mathbb{Z}^2}: t_\mathsf{L}=k, \mathsf{L}(0)=0 \} ] =\frac{4^{-k}}{k}\binom{k}{k/2}^2 \leq Ck^{-2}
    \end{equation}
    for a universal constant $C$.
    Let $\nu, C',C_1$ be the constants from Corollary~\ref{coro:vol}, let $\tilde C$ be the constant from Proposition \ref{prop:vol2} and let $n\geq 2$.
    For $R$ a non-negative integer, apply Corollary~\ref{coro:vol} to $k_R\coloneqq\lceil \max((R+1)\sqrt{k},  C' \log(n) )\rceil $:
    \begin{equation}
    \label{eq:temp:card}
    \mathbb{E}[\#\{ x\in n D: d(x,\gamma^{nD})\leq k_R  \}] \leq C_1  n^{2-\nu} k_R^{\nu}.
    \end{equation}
    For $\mathsf{L}$ rooted at $x$ to intersect $\gamma^{nD}$, it must hold $r_\mathsf{L}\geq d(x,\gamma^{nD})$. Let $\mathbb{E}^*$ denote the conditional expectation on the loop soup, given the path $\gamma^{nD}$ (i.e., we fix the realisation of $\gamma^{nD}$ and we take the expectation just over the loop soup $\mathbb{L}^{\mathbb{Z}^2}$. 
    By independence between $\mathbb{L}^{\mathbb{Z}^2}$ and $\gamma^{nD}$ and translation invariance of $\mathbb{L}^{\mathbb{Z}^2}$, for all $\delta\in(0,1/2]$, 
    \begin{align*}
    \mathbb{E}^*\Big[ &\#   \{   \mathsf{L}\in \mathbb{L}^{nD}_{\gamma^{nD}}: t_\mathsf{L}=k \}  \Big]
    = \sum_{R=0}^\infty
    \mathbb{E}^*\big[\# \{\mathsf{L}\in  \mathbb{L}^{nD}_{\gamma^{nD} } : t_\mathsf{L}=k, r_{\mathsf{L}}\in [R\sqrt{k}, (R+1)\sqrt{k}   )   \} \big]\\
    &	\leq \sum_{R=0}^\infty
    \mathbb{E}^*\big[\# \{\mathsf{L}\in  \mathbb{L}^{\mathbb{Z}^2} :
    \mathsf{L}(0)\in nD, \
    d(\mathsf{L}(0),\gamma^{nD})\leq (R+1)\sqrt{k}, \\
    &\hspace{6cm}
    t_\mathsf{L}=k, r_{\mathsf{L}}\in [R\sqrt{k}, (R+1)\sqrt{k}   )   \} \big]\\
    &	\leq \sum_{R=0}^\infty  \#
    \{  x\in n D : d(x,\gamma^{nD} )\leq  (R+1) \sqrt{k}  \}
    \\
    &\hspace{2cm}\cdot \mathbb{E}^*\big[\# \{L\in \mathbb{L}^{\mathbb{Z}^2}: L(0)=0, t_\mathsf{L}=k, r_{\mathsf{L}}\in [R\sqrt{k}, (R+1)\sqrt{k}  \} \big]\\
    &	\leq \sum_{R=0}^\infty  \#
    \big\{  x\in n D : d(x,\gamma^{nD} )\leq  (R+1) \sqrt{k}  \big\}
    \big(Ck^{-2} \big) \big( \tilde C \exp(-R^2) \big)\quad \text{by (\ref{eq:temp:mass},\ref{eq:temp:hoef})}
    \end{align*}
    Taking now the expectation with respect to $\gamma^{nD}$ and applying \eqref{eq:temp:card}, we get
    \begin{align*}
    \mathbb{E}\Big[ \#   \{   \mathsf{L}\in \mathbb{L}^{nD}_{\gamma^{nD}}: t_\mathsf{L}=k \}  \Big]
    % &	\leq \sum_{R=0}^\infty\mathbb{E}[ \#
    % \{  x\in n D : d(x,\gamma^{nD} )\leq  (R+1) \sqrt{k}  \}
    % ] \\
    % &\hspace{2cm} \cdot\mathbb{E}[\# \{L\in \mathbb{L}^{\mathbb{Z}^2}: L(0)=0, t_\mathsf{L}=k ]\cdot \mathbb{P}_k( r_{\mathsf{L}}\geq   R\sqrt{k}  ).\\
    &	\leq \sum_{R=0}^\infty
    \big( C_1  n^{2-\nu} k_R^{\nu} \big) \big(Ck^{-2} \big) 
    \big(\tilde C \exp(-R^2)\big) \\&=\tilde C CC_1 k^{-2}n^{2-\nu}  \sum_{R=0}^\infty  \exp(-R^2) k_R^{\nu}.
    \end{align*} 
    
    
    
    % \begin{align*}
    % \mathbb{E}\Big[ \#   \{   \mathsf{L}\in \mathbb{L}^{nD}_{\gamma^{nD}}: t_\mathsf{L}=k \}  \Big]
    % &	= \sum_{R=0}^\infty
    % \mathbb{E}\big[\# \{\mathsf{L}\in  \mathbb{L}^{nD}_{\gamma^{nD} } : t_\mathsf{L}=k, r_{\mathsf{L}}\in [R\sqrt{k}, (R+1)\sqrt{k}   )   \} \big]\\
    % &	\leq \sum_{R=0}^\infty
    % \mathbb{E}\big[\# \{\mathsf{L}\in  \mathbb{L}^{\mathbb{Z}^2} :
    % \mathsf{L}(0)\in nD, \
    % d(\mathsf{L}(0),\gamma^{nD})\leq (R+1)\sqrt{k}, \\
    % &\hspace{6cm}
    % t_\mathsf{L}=k, r_{\mathsf{L}}\in [R\sqrt{k}, (R+1)\sqrt{k}   )   \} \big]\\
    % &	\leq \sum_{R=0}^\infty \mathbb{E}\big[ \#
    % \{  x\in n D : d(x,\gamma^{nD} )\leq  (R+1) \sqrt{k}  \}
    % \big] \\
    % &\hspace{2cm}\cdot \mathbb{E}\big[\# \{L\in \mathbb{L}^{\mathbb{Z}^2}: L(0)=0, t_\mathsf{L}=k, r_{\mathsf{L}}\in [R\sqrt{k}, (R+1)\sqrt{k}  \} \big] \\
    % &	\leq \sum_{R=0}^\infty\mathbb{E}[ \#
    % \{  x\in n D : d(x,\gamma^{nD} )\leq  (R+1) \sqrt{k}  \}
    % ] \\
    % &\hspace{2cm} \cdot\mathbb{E}[\# \{L\in \mathbb{L}^{\mathbb{Z}^2}: L(0)=0, t_\mathsf{L}=k ]\cdot \mathbb{P}_k( r_{\mathsf{L}}\geq   R\sqrt{k}  ).\\
    % &	\leq \sum_{R=0}^\infty
    % \big( C_1  n^{2-\nu} k_R^{\nu} \big) \big(Ck^{-2} \big) \big( 4 \exp(-R^2) \big)\quad \text{by (\ref{eq:temp:card},\ref{eq:temp:mass},\ref{eq:temp:hoef})}\\
    % &=4CC_1 k^{-2}n^{2-\nu}  \sum_{R=0}^\infty  \exp(-R^2) k_R^{\nu}.
    % \end{align*}  %1+(R+1)\max(\sqrt{k},\log(n) )
    Since $k_R\leq 2 (R+1)(\max(\sqrt{k},\log(n) ))$ and $\sum_{R=0}^\infty   \exp(-R^2) (R+1)^\nu<\infty$, we deduce the first statement in Proposition \ref{prop:vol2}.
    
    Summing over $k=1$ to $\lceil \log(n)^2 \rceil$, and from $k=\lceil \log(n)^2 \rceil$ to $\delta n^2$ for an arbitrary $\delta>0$,
    we deduce that there exists a constant $C'$ such that for all $n\geq 1$ and all $\delta>0$,
    \[
    \mathbb{E}[T_{\leq \delta}(\mathcal{X}^n)]
    =n^{-2}
    \sum_{k=1}^{\lfloor \delta n^2 \rfloor}  k \mathbb{E}[\# \{  \mathsf{L}\in \mathbb{L}^{nD}_{\gamma^{nD}}: t_{\mathsf{L}}=k  \}]
    \leq C' ((\log\log n) (\log n)^\nu n^{-\nu}+ \delta^{\frac{\nu}{2}}).
    %\qedhere
    \]
    This proves \eqref{eq:temp:Ebound}.
\end{proof}

\end{document}
